% ==================================================================
% id:       func-005
% topic:    פונקציות
% subtopic: תחום הגדרה ותכונות של פונקציות
% solution: full
% ==================================================================
\textbf{קבעו מהו תחום ההגדרה הטבעי של הפונקציות הבאות:}
\begin{enumerate}\bMs{2}
\item  $f(x) = (-x)^{\frac{3}{2}}$
\begin{solbox}
$(- \infty, 0]$
\end{solbox}
\item  $f(x) = \sqrt{\sin(\sqrt{x})}$
\hint{דרשו ש-$x\geq0$ וש-$\sin(\sqrt x)\geq0$. פתרו קודם עבור $t=\sqrt x$ ואחר כך העלו בריבוע.}
\begin{solbox}
$\bigcup_{k=0}^{\infty} [4k^2 \pi^2, (2k + 1)^2 \pi^2 ]$
\end{solbox}
\item  $f(x) = \sqrt[4]{\log_2(\tan x)}$
\hint{דרשו שהביטוי שבתוך הלוגריתם יהיה חיובי ושהלוגריתם עצמו יהיה אי-שלילי, והיעזרו במחזוריות של $\tan$.}
\begin{solbox}
ממחזוריות של
$\tan$
נוכל לפתור את השאלה בקטע
$(-\frac{\pi}{2}, -\frac{\pi}{2})$
ואז לשכפל את התוצאה על הישר כולו.
כדי שפונקציית הלוג תוגדר, היא צריכה לקבל ערכים חיוביים. לכן נדרוש
$\tan x \geq 0$.
זה קורה בקטע
$(0, \frac{\pi}{2})$.
כדי שנוכל להוציא שורש רביעי ממספר הוא צריך להיות אי שלילי. לכן נדרוש
$\log_2 (\tan x) \geq 0$.
זה קורה כאשר
$\tan x \geq 1$,
וזה קורה בקטע
$[\frac{\pi}{4}, \frac{\pi}{2})$.
כעת נשכפל את מה שקיבלנו לישר, ונקבל שתחום ההגדרה הוא
  $\bigcup_{k=-\infty}^{\infty} [\frac{\pi}{4} + k \pi, \frac{\pi}{2} + k \pi )$
\end{solbox}
\eMs
\end{enumerate}
